Video and world models can imagine futures that look plausible and reach a goal, yet require local changes the robot body cannot produce. We study this failure mode through \emph{actionability energy}: the distance between an imagined representation transition and the action-reachable tangent set induced by an embodiment action map. The central claim is that Jacobian-style action maps, such as those used for visual inverse control, can also serve as diagnostic interfaces for world-model planning. They can certify whether a future is locally executable, predict failure, and repair non-actionable trajectories by compositional inference. We implement analytic and learned local Jacobian maps on a CPU benchmark with a point mass, nonholonomic car, two-link arm, contact-switching pusher, and rendered keypoint futures. Analytic actionability is near-perfect, learned local action maps recover most of the diagnostic signal, and a learned classifier can fit the synthetic benchmark but does not explain the mechanism. Smoothness, passive likelihood, action norm, inverse-dynamics reconstruction, wrong-map residuals, and forward-inverse asymmetry proxy baselines are weaker. Robustness audits expose the hard case: discontinuous contact modes. The study is intentionally controlled and does not claim real-robot validation or large video-model finetuning; it isolates a mechanism for making generated futures honest about the body that must execute them.
